\documentclass[11pt]{article}
\usepackage[a4paper,margin=22mm]{geometry}
\usepackage{fontspec}
\setmainfont{DejaVu Serif}
\setsansfont{DejaVu Sans}
\usepackage{microtype}
\usepackage{xcolor}
\usepackage{hyperref}
\usepackage{enumitem}
\usepackage{parskip}
\definecolor{Accent}{HTML}{971D32}
\definecolor{Muted}{HTML}{666666}
\hypersetup{colorlinks=true,urlcolor=Accent,linkcolor=Accent}
\setlist{leftmargin=1.6em,itemsep=0.3em,topsep=0.35em}
\setcounter{tocdepth}{2}
\renewcommand{\contentsname}{Contents}
\newcommand{\sectionrule}{\par\vspace{-0.5em}\noindent\textcolor{Accent}{\rule{\linewidth}{0.6pt}}\par}
\begin{document}

\begin{center}
{\sffamily\bfseries\color{Accent} TOEFL WORD OF THE DAY\par}
\vspace{8mm}
{\fontsize{42}{48}\selectfont\bfseries complacent\par}
\vspace{2mm}
{\Large\sffamily\color{Accent} /kəmˈpleɪsənt/\par}
\vspace{2mm}
{\small\sffamily Local audio phrase: complacent\par}
{\small\sffamily Audio file: audios/complacent.mp3\par}
\vspace{2mm}
{\large\sffamily\color{Muted} adjective\par}
\vspace{5mm}
{\sffamily too satisfied to notice risks or improve\par}
\vspace{6mm}
{\large A complacent person feels so satisfied or secure that they stop paying attention to problems, risks, or the need for improvement.\par}
\vspace{6mm}
\begin{minipage}{0.84\textwidth}
\itshape “After several successful years, the company became complacent and stopped investing in innovation.”
\end{minipage}\par
\vspace{10mm}
\textbf{Date:} August 6th 2026\quad\textbullet\quad \textbf{Level:} B1--mid-B2 TOEFL
\vfill
Created and compiled by \textbf{Amir Kooshky}\par
\vspace{2mm}
\href{https://t.me/KooshkyTOEFL}{Telegram Channel}\quad\textbullet\quad
\href{https://www.instagram.com/kooshkytoefl}{Instagram}\quad\textbullet\quad
\href{https://t.me/KooshkyTOEFL\_pv}{Personal Telegram}
\end{center}
\newpage
\tableofcontents
\newpage

\section{Meanings and usage}
\sectionrule
\subsection{Meaning 1: Too satisfied to notice risks}
\textbf{Part of speech:} adjective

\textbf{IPA:} /kəmˈpleɪsənt/

\textbf{Pronunciation phrase:} complacent

\textbf{Local audio file:} audios/complacent.mp3

\textbf{Register:} neutral to formal; common in academic, professional, political, health, and safety discussions

\textbf{Definition:} too satisfied with a situation or with your own abilities, so that you do not notice problems, make improvements, or prepare for possible danger

\textbf{In simpler words:} too comfortable and no longer careful

\textbf{Typical pattern:} become or grow complacent about + problem, risk, or situation

\subsubsection{Natural examples}
\begin{enumerate}
  \item After several successful years, the company became complacent and stopped investing in innovation.
  \item Public health officials warned people not to become complacent about the disease.
  \item The team grew complacent after taking an early lead and eventually lost the game.
  \item Students should not become complacent simply because they performed well on one examination.
  \item The report argues that economic growth made policymakers complacent about rising inequality.
  \item Strong sales may create a complacent attitude toward changes in consumer behavior.
  \item The absence of recent accidents should not make workers complacent about safety.
  \item Competitors quickly gained ground while the market leader remained complacent.
\end{enumerate}

\subsubsection{Nuance and implication}
\textbf{Central nuance:} Complacent does not simply mean calm, comfortable, or satisfied. It means that satisfaction has become excessive and has reduced attention, effort, or caution.

\textbf{Usual implication:} The word is negative and usually suggests that the person or group may fail, fall behind, or face preventable problems.

\textbf{Compared with content:} Content means peacefully satisfied and is usually positive or neutral. Complacent means too satisfied to recognize risks or continue improving and is negative.

\subsubsection{Grammar notes}
\textbf{How to use it:} Use it after a linking verb such as be, become, grow, seem, or remain, or before a noun such as attitude or leadership.

\textbf{What can follow it:} It is often followed by about and then a problem, risk, achievement, or situation.

\textbf{Common preposition:} Use complacent about something, as in complacent about safety. Do not normally use complacent of or complacent with.

\textbf{Position in a sentence:} It commonly appears after become or grow when describing a change in attitude.

\subsubsection{Useful patterns}
\textbf{Pattern:} become or grow complacent

\textbf{Use:} Use this when someone gradually becomes less careful or ambitious because they feel too secure.

\textbf{Example:} The organization grew complacent after years without serious competition.

\textbf{Pattern:} be complacent about + noun

\textbf{Use:} Use this to identify the problem, risk, or situation that someone is not taking seriously enough.

\textbf{Example:} We cannot afford to be complacent about cybersecurity.

\textbf{Pattern:} be complacent about + -ing form

\textbf{Use:} Use this when describing an activity or development that someone is treating with too little concern.

\textbf{Example:} Officials were complacent about allowing the infrastructure to deteriorate.

\textbf{Pattern:} a complacent attitude

\textbf{Use:} Use this for a way of thinking that ignores danger or the need for improvement.

\textbf{Example:} A complacent attitude toward workplace safety can have serious consequences.

\textbf{Pattern:} remain complacent despite + noun

\textbf{Use:} Use this when someone continues to feel too secure even though warning signs exist.

\textbf{Example:} The company remained complacent despite a steady decline in customer satisfaction.

\subsubsection{Common wrong usage}
\paragraph{Correction 1}
\textbf{Wrong:} The students were complacent with the coming examination.

\textbf{Correct:} The students were complacent about the coming examination.

\textbf{Why:} Use complacent about to identify the matter that someone is not taking seriously enough.

\paragraph{Correction 2}
\textbf{Wrong:} She felt complacent because she had completed a difficult project and deserved to relax.

\textbf{Correct:} She felt satisfied because she had completed a difficult project and deserved to relax.

\textbf{Why:} Complacent is negative and suggests a harmful lack of effort or caution. Use satisfied when the feeling is reasonable.

\paragraph{Correction 3}
\textbf{Wrong:} The manager complacented after the company became successful.

\textbf{Correct:} The manager became complacent after the company became successful.

\textbf{Why:} Complacent is an adjective, not a verb. Use become complacent.

\paragraph{Correction 4}
\textbf{Wrong:} The government should not be complacent of the growing risk.

\textbf{Correct:} The government should not be complacent about the growing risk.

\textbf{Why:} Use about, not of, after complacent.

\section{Word family}
\sectionrule
\subsection{complacency}
\textbf{Part of speech:} uncountable noun

\textbf{IPA:} /kəmˈpleɪsənsi/

\textbf{Definition:} a feeling of excessive satisfaction that causes someone to ignore problems, risks, or the need for improvement

\textbf{Relationship to the headword:} This noun names the harmful attitude or state described by complacent.

\textbf{Grammar pattern:} complacency about + issue or a sense of complacency

\textbf{Usage note:} It is normally uncountable, so do not usually use a or an before it.

\subsubsection{Examples}
\begin{enumerate}
  \item The recent improvement should not lead to complacency.
  \item Officials warned that public complacency could allow the disease to spread again.
  \item The company's decline was partly caused by years of managerial complacency.
\end{enumerate}

\subsection{complacently}
\textbf{Part of speech:} adverb

\textbf{IPA:} /kəmˈpleɪsəntli/

\textbf{Definition:} in a way that shows excessive satisfaction and too little awareness of problems or risks

\textbf{Relationship to the headword:} This adverb describes an action performed with a complacent attitude.

\textbf{Grammar pattern:} complacently + assume, believe, ignore, or observe

\textbf{Usage note:} It is less common than complacent and complacency.

\subsubsection{Examples}
\begin{enumerate}
  \item The company complacently assumed that customers would remain loyal.
  \item Officials watched complacently as the warning signs became more serious.
\end{enumerate}

\section{Common collocations}
\sectionrule
\subsection{become complacent}
\textbf{Applies to:} Too satisfied to notice risks

\textbf{Meaning:} begin to feel so secure that effort or caution decreases

\textbf{Pattern:} become complacent about + issue

\textbf{Example:} Successful organizations can fail if they become complacent.

\subsection{grow complacent}
\textbf{Applies to:} Too satisfied to notice risks

\textbf{Meaning:} gradually become too satisfied and less careful

\textbf{Example:} The team grew complacent after winning several easy matches.

\subsection{complacent about}
\textbf{Applies to:} Too satisfied to notice risks

\textbf{Meaning:} not sufficiently concerned or careful about something

\textbf{Pattern:} be complacent about + noun or -ing form

\textbf{Example:} Consumers should not be complacent about online privacy.

\subsection{a complacent attitude}
\textbf{Applies to:} Too satisfied to notice risks

\textbf{Meaning:} a way of thinking that ignores danger or the need for improvement

\textbf{Example:} A complacent attitude toward maintenance increased the risk of equipment failure.

\subsection{complacent leadership}
\textbf{Applies to:} Too satisfied to notice risks

\textbf{Meaning:} leaders who are too satisfied with current conditions to recognize problems

\textbf{Example:} Complacent leadership prevented the organization from responding quickly to change.

\subsection{a sense of complacency}
\textbf{Applies to:} Too satisfied to notice risks

\textbf{Meaning:} a feeling that there is no need to worry or make further effort

\textbf{Example:} The long period of stability created a false sense of complacency.

\textbf{Usage note:} Complacency is the noun form.

\subsection{lead to complacency}
\textbf{Applies to:} Too satisfied to notice risks

\textbf{Meaning:} cause people to become excessively satisfied and less careful

\textbf{Pattern:} success, improvement, or stability + lead to complacency

\textbf{Example:} Early success can lead to complacency if progress is not carefully monitored.

\subsection{guard against complacency}
\textbf{Applies to:} Too satisfied to notice risks

\textbf{Meaning:} take action to prevent excessive confidence or lack of caution

\textbf{Pattern:} guard against complacency

\textbf{Example:} Regular safety training helps workers guard against complacency.

\subsection{room for complacency}
\textbf{Applies to:} Too satisfied to notice risks

\textbf{Meaning:} an acceptable reason or opportunity to become less careful

\textbf{Pattern:} there is no room for complacency

\textbf{Example:} Although infection rates have fallen, there is no room for complacency.

\subsection{avoid complacency}
\textbf{Applies to:} Too satisfied to notice risks

\textbf{Meaning:} prevent an attitude of excessive satisfaction

\textbf{Example:} The industry must continue to innovate to avoid complacency.

\section{Synonyms and antonyms}
\sectionrule
\subsection{Close synonyms}
\subsubsection{self-satisfied}
\textbf{Part of speech:} adjective

\textbf{Applies to:} Too satisfied to notice risks

\textbf{Shared meaning:} Both describe an excessive and usually negative form of satisfaction.

\textbf{Difference:} Self-satisfied emphasizes excessive satisfaction with yourself or your achievements. Complacent also emphasizes the resulting failure to notice risks or continue improving.

\textbf{Example:} His self-satisfied expression annoyed the other members of the team.

\subsubsection{smug}
\textbf{Part of speech:} adjective

\textbf{Applies to:} Too satisfied to notice risks

\textbf{Shared meaning:} Both can describe someone who is excessively pleased with themselves.

\textbf{Difference:} Smug strongly suggests that someone feels superior and often shows it in an irritating way. Complacent focuses more on reduced caution, effort, or awareness.

\textbf{Example:} She sounded smug after her prediction proved correct.

\subsubsection{overconfident}
\textbf{Part of speech:} adjective

\textbf{Applies to:} Too satisfied to notice risks

\textbf{Shared meaning:} Both attitudes can cause someone to underestimate difficulties or risks.

\textbf{Difference:} Overconfident means having more confidence than the situation justifies. Complacent means feeling so satisfied or secure that you stop noticing problems or making an effort.

\textbf{Example:} The candidate became overconfident and failed to prepare for the final interview.

\subsection{Direct or useful antonyms}
\subsubsection{vigilant}
\textbf{Part of speech:} adjective

\textbf{Applies to:} Too satisfied to notice risks

\textbf{Opposition:} A vigilant person remains watchful for possible problems or danger, while a complacent person fails to give them enough attention.

\textbf{Example:} Health officials remained vigilant for signs of another outbreak.

\subsubsection{diligent}
\textbf{Part of speech:} adjective

\textbf{Applies to:} Too satisfied to notice risks

\textbf{Opposition:} A diligent person continues to work carefully and consistently, while a complacent person may reduce their effort because they feel too secure.

\textbf{Usage note:} This is an opposite mainly when complacent refers to a lack of continued effort.

\textbf{Example:} The research team was diligent in checking every stage of the experiment.

\section{Words learners may confuse}
\sectionrule
\subsection{complaisant}
\textbf{Part of speech:} adjective

\textbf{Why learners confuse them:} The two words have similar spellings and pronunciations but different meanings.

\textbf{Key difference:} Complacent means too satisfied and therefore insufficiently careful. Complaisant means willing to please other people or agree with them.

\textbf{complacent:} The company became complacent after dominating the market for years.

\textbf{complaisant:} The complaisant assistant agreed to every request.

\subsection{compliant}
\textbf{Part of speech:} adjective

\textbf{Why learners confuse them:} Both words begin with compl- and may occur in formal professional contexts.

\textbf{Key difference:} Complacent describes excessive satisfaction and lack of caution. Compliant describes someone or something that follows rules, requests, or standards.

\textbf{complacent:} Managers became complacent about the growing security threat.

\textbf{compliant:} The new equipment is fully compliant with safety regulations.

\section{Practice exercises}
\sectionrule
\subsection{Word family choice}
Choose the best member of the word family for each blank.

\begin{enumerate}
  \item The recent decline in accidents must not create a sense of \_\_\_\_\_.
  \item The organization became \_\_\_\_\_ after several years of rapid growth.
  \item Executives \_\_\_\_\_ assumed that no competitor could challenge the company.
\end{enumerate}

\subsection{Correct or incorrect?}
Decide whether each sentence is correct. Correct the inaccurate ones.

\begin{enumerate}
  \item We should not become complacent about workplace safety.
  \item The company was complacent with changes in consumer demand.
  \item She was complacent because she was reasonably pleased with her progress.
  \item The absence of immediate danger led to complacency.
\end{enumerate}

\subsection{Collocation practice}
Complete each sentence with a natural collocation.

\begin{enumerate}
  \item A long period without accidents can lead to \_\_\_\_\_.
  \item The public must not become complacent \_\_\_\_\_ the risk of flooding.
  \item Regular inspections help guard \_\_\_\_\_ complacency.
  \item Despite recent improvements, there is no room for \_\_\_\_\_.
\end{enumerate}

\subsection{Complete answer key}
\subsubsection{Word family choice}
\begin{enumerate}
  \item \textbf{complacency.} The recent decline in accidents must not create a sense of complacency. \emph{Use the noun complacency after a sense of.}
  \item \textbf{complacent.} The organization became complacent after several years of rapid growth. \emph{Use the adjective complacent after became.}
  \item \textbf{complacently.} Executives complacently assumed that no competitor could challenge the company. \emph{Use the adverb complacently to describe how they assumed something.}
\end{enumerate}

\subsubsection{Correct or incorrect?}
\begin{enumerate}
  \item \textbf{Correct} \emph{Become complacent about is the standard pattern.}
  \item \textbf{Incorrect}. The company was complacent about changes in consumer demand. \emph{Use complacent about, not complacent with.}
  \item \textbf{Incorrect}. She was satisfied because she was reasonably pleased with her progress. \emph{Complacent is negative and suggests harmful overconfidence or lack of effort.}
  \item \textbf{Correct} \emph{Lead to complacency is a natural expression for causing excessive confidence or reduced caution.}
\end{enumerate}

\subsubsection{Collocation practice}
\begin{enumerate}
  \item \textbf{complacency.} A long period without accidents can lead to complacency. \emph{Lead to complacency means cause people to become less careful because they feel secure.}
  \item \textbf{about.} The public must not become complacent about the risk of flooding. \emph{Use complacent about something.}
  \item \textbf{against.} Regular inspections help guard against complacency. \emph{Guard against complacency means take action to prevent it.}
  \item \textbf{complacency.} Despite recent improvements, there is no room for complacency. \emph{There is no room for complacency means people must remain careful and continue making an effort.}
\end{enumerate}

\vfill
\begin{center}
\small\color{Muted} Created and compiled by \textbf{Amir Kooshky}\\
\href{https://t.me/KooshkyTOEFL}{Telegram Channel}\quad\textbullet\quad
\href{https://www.instagram.com/kooshkytoefl}{Instagram}\quad\textbullet\quad
\href{https://t.me/KooshkyTOEFL\_pv}{Personal Telegram}
\end{center}
\end{document}
